\documentclass[12pt, twoside]{article}
\input{../preamble}

\fancyhead[LE]{\thepage}
\fancyhead[RO]{Markscheme}
\fancyhead[LO]{La Scuola d'Italia / Dr.\ Huson / IB Math \\* 16 April 2026}

\begin{document}
\subsubsection*{6.7 PreQuiz: Polynomials --- Markscheme}

\begin{enumerate}[itemsep=0.3cm]

%--- Problem 1: f(x) = x^2 + 2x - 15 [7 marks] ---
\item[1.] [Maximum mark: 7]

Let $f(x) = x^2 + 2x - 15$.

\medskip

\textbf{(a)} Write $f(x)$ in factored form $(x - a)(x - b)$. \hfill [2]

$f(x) = (x - 3)(x - (-5)) = (x - 3)(x + 5)$ \hfill \textbf{(M1)(A1)}

\emph{Accept $(x + 5)(x - 3)$.}

\medskip

\textbf{(b)} Write down the $x$-intercepts. \hfill [1]

$x = 3$ and $x = -5$ \hfill \textbf{(A1)}

\medskip

\textbf{(c)} Write down the $y$-intercept. \hfill [1]

$f(0) = -15$, so the $y$-intercept is $(0, -15)$. \hfill \textbf{(A1)}

\medskip

\textbf{(d)} Write $f(x)$ in vertex form $(x - h)^2 + k$. \hfill [1]

$f(x) = (x + 1)^2 - 16$ \hfill \textbf{(A1)}

\emph{Or equivalently, $(x - (-1))^2 + (-16)$, i.e.\ $h = -1$, $k = -16$.}

\medskip

\textbf{(e)} Find the coordinates of the vertex. \hfill [2]

Axis of symmetry: $x = \dfrac{-5 + 3}{2} = -1$ \hfill \textbf{(A1)}

$f(-1) = 1 - 2 - 15 = -16$

Vertex: $(-1, -16)$ \hfill \textbf{(A1)}

\bigskip
\hrule
\bigskip

%--- Problem 2: g(x) = -(x - 4)^2 + 9 [7 marks] ---
\item[2.] [Maximum mark: 7]

Consider the function $g(x) = -(x - 4)^2 + 9$.

\medskip

\textbf{(a)} Write down the coordinates of the vertex. \hfill [1]

Vertex: $(4, 9)$ \hfill \textbf{(A1)}

\medskip

\textbf{(b)} Write down the equation of the axis of symmetry. \hfill [1]

$x = 4$ \hfill \textbf{(A1)}

\medskip

\textbf{(c)} Show that $g(x) = -x^2 + 8x - 7$. \hfill [2]

$g(x) = -(x - 4)^2 + 9 = -(x^2 - 8x + 16) + 9$ \hfill \textbf{(M1)}

$= -x^2 + 8x - 16 + 9 = -x^2 + 8x - 7$ \hfill \textbf{(A1)(AG)}

\medskip

\textbf{(d)} Solve $g(x) = 0$ to find the $x$-intercepts. \hfill [3]

$-x^2 + 8x - 7 = 0$ \hfill \textbf{(M1)}

$x^2 - 8x + 7 = 0$

$(x - 1)(x - 7) = 0$ \hfill \textbf{(A1)}

$x = 1$ or $x = 7$ \hfill \textbf{(A1)}

\newpage

%--- Problem 3: y = 2x and y = x^2 - 3 [7 marks] ---
\item[3.] [Maximum mark: 7]

The line $y = 2x$ and the parabola $y = x^2 - 3$.

\medskip

\textbf{(a)} Show that the $x$-coordinates of intersection
satisfy $x^2 - 2x - 3 = 0$. \hfill [2]

At intersection: $x^2 - 3 = 2x$ \hfill \textbf{(M1)}

$x^2 - 2x - 3 = 0$ \hfill \textbf{(A1)(AG)}

\medskip

\textbf{(b)} Show that $x^2 - 2x - 3 = 0$ has two distinct
real roots. \hfill [1]

$\Delta = (-2)^2 - 4(1)(-3) = 4 + 12 = 16 > 0$ \hfill \textbf{(A1)}

Since $\Delta > 0$, there are two distinct real roots.

\medskip

\textbf{(c)} Factorise $x^2 - 2x - 3$. \hfill [1]

$(x - 3)(x + 1)$ \hfill \textbf{(A1)}

\medskip

\textbf{(d)} Find the $x$-coordinates of the points
of intersection. \hfill [1]

$x = 3$ and $x = -1$ \hfill \textbf{(A1)}

\medskip

\textbf{(e)} Find the coordinates of both points
of intersection. \hfill [2]

$x = 3$: $y = 2(3) = 6 \implies (3, 6)$ \hfill \textbf{(A1)}

$x = -1$: $y = 2(-1) = -2 \implies (-1, -2)$ \hfill \textbf{(A1)}

\bigskip
\hrule
\bigskip

%--- Problem 4: 2x^2 - kx + 8 = 0 [7 marks] ---
\item[4.] [Maximum mark: 7]

The equation $2x^2 - kx + 8 = 0$ has two equal roots.

\medskip

\textbf{(a)} Write down the discriminant in terms of $k$. \hfill [1]

$\Delta = (-k)^2 - 4(2)(8) = k^2 - 64$ \hfill \textbf{(A1)}

\medskip

\textbf{(b)} Find the values of $k$. \hfill [2]

Equal roots: $\Delta = 0$

$k^2 - 64 = 0 \implies k^2 = 64$ \hfill \textbf{(M1)}

$k = \pm 8$ \hfill \textbf{(A1)}

\medskip

\textbf{(c)} For $k = 8$, find the repeated root. \hfill [2]

$2x^2 - 8x + 8 = 0 \implies x^2 - 4x + 4 = 0$ \hfill \textbf{(M1)}

$(x - 2)^2 = 0 \implies x = 2$ \hfill \textbf{(A1)}

\medskip

\textbf{(d)} State the values of $k$ for which the equation
has no real roots. \hfill [2]

$\Delta < 0 \implies k^2 - 64 < 0 \implies k^2 < 64$ \hfill \textbf{(M1)}

$-8 < k < 8$ \hfill \textbf{(A1)}

\end{enumerate}
\end{document}
